\begin{itemize}
    \item \textbf{Lucas's Theorem} \\
    \begin{equation*}
    \begin{aligned}
    \binom{n}{k} \equiv \prod_{i=0}^{r} \binom{n_i}{k_i} \pmod{p}
    \end{aligned}
    \end{equation*} \\
    where $p$ is the prime and $n=n_m n_{m-1} \cdots n_0$ and $k=k_m k_{m-1} \cdots k_0$ are the base representations of $p$.

    \item \textbf{Euler's Theorem} \\
    \begin{equation*}
    \begin{aligned}
    a^{\varphi(n)} \equiv 1 \pmod{n} \quad \text{for } \gcd(a, n) = 1
    \end{aligned}
    \end{equation*}

    \item \textbf{Möbius Inversion Formula} \\ 
    \begin{equation*}
    \begin{aligned}
    f(n) = \sum_{d \mid n} g(d)
\quad \Longleftrightarrow \quad
g(n) = \sum_{d \mid n} \mu(d)\, f\!\left(\frac{n}{d}\right)
    \end{aligned}
    \end{equation*}
    \begin{equation*}
    \begin{aligned}
    \sum_{d|n}\phi(d) = n \quad \Longleftrightarrow \quad \phi(n)=\sum_{d|n}\mu(d)\frac{n}{d}
    \end{aligned}
    \end{equation*}

    \item \textbf{Green's Theorem} \\
    \begin{equation*}
    \begin{aligned}
    \oint_{C} \left( P\,dx + Q\,dy \right)
=
\iint_{D} \left( \frac{\partial Q}{\partial x} - \frac{\partial P}{\partial y} \right)\, dA
    \end{aligned}
    \end{equation*}

    \item \textbf{Fermat's Little Theorem} \\
    \begin{equation*}
    \begin{aligned}
    a^{p-1} \equiv 1 \pmod{p} \quad \text{for } \gcd(a, p) = 1
    \end{aligned}
    \end{equation*} \\ 
    Equivalently, $a^p \equiv a \pmod{p}$ for all integers $a$.

    \item \textbf{Pick's theorem} \\
    \begin{equation*}
    \begin{aligned}
    A = I + \frac{B}{2} - 1
    \end{aligned}
    \end{equation*} \\ 
    where $A$ is the area of a simple polygon whose vertices are points on a lattice. $I$ is the number of interior lattice points, and $B$ is the number of boundary lattice points. 

    \item \textbf{Burnside's Lemma} \\
    \begin{equation*}
    \begin{aligned}
    |X/G| = \frac{1}{|G|} \sum_{g \in G} |X^g|
    \end{aligned}
    \end{equation*} \\ 
    where $X$ is a set, $G$ is a group acting on $X$. $X/G$ is the set of orbits, and $X^g$ is the set of elements fixed by $g$.

    \item \textbf{Kirchhoff's Theorem} \\
    \begin{equation*}
    \begin{aligned}
    \text{Number of spanning trees } = \det(L^*)
    \end{aligned}
    \end{equation*} \\ 
    where $L$ is the Laplacian matrix of the graph and $L^*$ is any cofactor of $L$.

    \item \textbf{Cayley's Formula} \\
    \begin{equation*}
    \begin{aligned}
    \text{Number of labeled trees with $n$ vertices} = n^{n-2}
    \end{aligned}
    \end{equation*}

    \item \textbf{Simpson's Rule} \\
    \begin{equation*}
    \begin{aligned}
    \int_{a}^{b} f(x)\,dx \approx \frac{b-a}{6}
\left[ f(a) + 4f\!\left(\frac{a+b}{2}\right) + f(b) \right]
    \end{aligned}
    \end{equation*} 

    \item \textbf{Law of Sines} \\
    \begin{equation*}
    \begin{aligned}
    \frac{a}{\sin A} = \frac{b}{\sin B} = \frac{c}{\sin C}
    \end{aligned}
    \end{equation*}

    \item \textbf{Law of Cosines} \\ 
    \begin{equation*}
    \begin{aligned}
    c^2 = a^2 + b^2 - 2ab \cos C
    \end{aligned}
    \end{equation*}

    \item \textbf{Area (Heron's Formula)} \\ 
    \begin{equation*}
    \begin{aligned}
    A = \sqrt{s(s-a)(s-b)(s-c)} \quad \text{where } s = \frac{a+b+c}{2}
    \end{aligned}
    \end{equation*}

    \item \textbf{Point and Line Distance} \\ 
    \begin{equation*}
    \begin{aligned}
    d = \frac{|Ax_0 + By_0 + C|}{\sqrt{A^2 + B^2}}
    \end{aligned}
    \end{equation*}

    \item \textbf{Angle Addition Formulas} \\
    \begin{equation*}
    \begin{aligned}
    \sin(A \pm B) &= \sin A \cos B \pm \cos A \sin B \\
\cos(A \pm B) &= \cos A \cos B \mp \sin A \sin B \\
\tan(A \pm B) &= \frac{\tan A \pm \tan B}{1 \mp \tan A \tan B}
    \end{aligned}
    \end{equation*}

    \item \textbf{Double Angle Formulas} \\ 
    \begin{equation*}
    \begin{aligned}
    \sin(2x) &= 2\sin x \cos x \\
\cos(2x) &= \cos^2 x - \sin^2 x = 2\cos^2 x - 1 = 1 - 2\sin^2 x \\
\tan(2x) &= \frac{2\tan x}{1 - \tan^2 x}
    \end{aligned}
    \end{equation*} 

    \item \textbf{Half Angle Formulas} \\ 
    \begin{equation*}
    \begin{aligned}
    \sin\left(\frac{x}{2}\right) = \pm \sqrt{\frac{1 - \cos x}{2}}, \quad
\cos\left(\frac{x}{2}\right) = \pm \sqrt{\frac{1 + \cos x}{2}}
    \end{aligned}
    \end{equation*}
    \begin{equation*}
    \begin{aligned}
    \tan\left(\frac{x}{2}\right) = \frac{\sin x}{1 + \cos x}
    \end{aligned}
    \end{equation*}

    \item \textbf{Catalan Numbers} \\ 
    \begin{equation*}
    \begin{aligned}
    C_n = \frac{1}{n+1}\binom{2n}{n}
    \end{aligned}
    \end{equation*}
    \begin{equation*}
    \begin{aligned}
    C_0 = 1, \quad C_{n+1}=\sum_{i=0}^n C_i C_{n-i} \quad \text{for } n \geq 0
    \end{aligned}
    \end{equation*} \\ 
    $1, 1, 2, 5, 14, 42, 132, 429, 1430, 4862, 16796, 58786$

    \item \textbf{Ballot Number} \\
    \begin{equation*}
    \begin{aligned}
    B(p, q) = \frac{p-q}{p+q}\binom{p+q}{p} \quad (p > q \geq 0)
    \end{aligned}
    \end{equation*} \\
    counts vote sequences where A stays strictly ahead throughout; $B(n+1,n) = C_n$.

    \item \textbf{Motzkin Numbers}
\begin{minted}[fontsize=\small]{c++}
m[0] = 1;
for(int n=0; n<N; n++){
    m[n+1] = m[n];
    for(int k=0; k<n; k++) m[n+1] = (m[n+1] + m[k] * m[n-1-k]) % MOD;
}
\end{minted}
    counts paths $(0,0) \to (n,0)$ with steps $U=(1,1), D=(1,-1), F=(1,0)$ that never go below the $x$-axis.

    \item \textbf{Narayana Numbers} \\
    \begin{equation*}
    \begin{aligned}
    N(n, k) = \frac{1}{n}\binom{n}{k}\binom{n}{k-1} \quad (1 \leq k \leq n)
    \end{aligned}
    \end{equation*} \\
    counts Dyck paths of length $2n$ with exactly $k$ peaks; $\sum_{k=1}^n N(n,k) = C_n$.

    \item \textbf{Schr\"oder Numbers (large)} \\
    \begin{equation*}
    \begin{aligned}
    S_0 = 1, \ S_1 = 2, \ (n+1)S_n = 3(2n-1)S_{n-1} - (n-2)S_{n-2} \quad (n \geq 2)
    \end{aligned}
    \end{equation*} \\
    counts monotone paths $(0,0) \to (n,n)$ with steps $(1,0), (0,1), (1,1)$ that never rise above $y=x$.

    \item \textbf{Stirling Numbers (1st kind, unsigned)}
\begin{minted}[fontsize=\small]{c++}
s[0][0] = 1;
for(int n=1; n<=N; n++) for(int k=1; k<=n; k++)
    s[n][k] = ((n-1) * s[n-1][k] + s[n-1][k-1]) % MOD;
\end{minted}
    unsigned $\left[{n\atop k}\right]$ counts permutations of $n$ elements with $k$ cycles.

    \item \textbf{Stirling Numbers (2nd kind)}
\begin{minted}[fontsize=\small]{c++}
s[0][0] = 1;
for(int n=1; n<=N; n++) for(int k=1; k<=n; k++)
    s[n][k] = (k * s[n-1][k] + s[n-1][k-1]) % MOD;
\end{minted}
    counts partitions of an $n$-set into $k$ nonempty unlabeled blocks.

    \item \textbf{Eulerian Numbers}
\begin{minted}[fontsize=\small]{c++}
a[0][0] = 1;
for(int n=1; n<=N; n++) for(int k=0; k<n; k++)
    a[n][k] = ((k+1) * a[n-1][k] + (k ? (n-k) * a[n-1][k-1] : 0)) % MOD;
\end{minted}
    counts permutations of $1..n$ with exactly $k$ descents ($0 \leq k \leq n-1$).

    \item \textbf{Bell Numbers} \\
    \begin{equation*}
    \begin{aligned}
    B_0 = 1, \quad B_{n+1} = \sum_{k=0}^n \binom{n}{k} B_k = \sum_{k=0}^n \left\{{n\atop k}\right\}
    \end{aligned}
    \end{equation*} \\
    counts partitions of an $n$-set into any number of nonempty unlabeled blocks.
    Equivalently, $B_n = n! \cdot [x^n]\, e^{e^x - 1}$: expand $e^{e^x-1}$ as a formal power series and take $n!$ times the coefficient of $x^n$.

    \item \textbf{Partition Numbers}
\begin{minted}[fontsize=\small]{c++}
p[0] = 1;
for(int n=1; n<=N; n++) for(int k=1, g; (g = k*(3*k-1)/2)<=n; k++){
    ll t = (p[n-g] + (g+k <= n ? p[n-g-k] : 0)) % MOD;
    p[n] = (p[n] + (k & 1 ? t : MOD - t)) % MOD;
}
\end{minted}
    (pentagonal number theorem) counts ways to write $n$ as a sum of positive integers, order not mattering.

    \item \textbf{Parking Functions} \\
    \begin{equation*}
    \begin{aligned}
    \text{count of length-}n\text{ parking functions} = (n+1)^{n-1}
    \end{aligned}
    \end{equation*} \\
    $(a_i)$, $1 \leq a_i \leq n$, is one iff sorted $b_1 \leq \cdots \leq b_n$ has $b_i \leq i$ for all $i$.

    \item \textbf{Fibonacci Numbers} \\
    \begin{equation*}
    \begin{aligned}
    F_0 = 0, \quad F_1 = 1, \quad F_n=F_{n-1} + F_{n-2} \quad \text{for } n \geq 2 
    \end{aligned}
    \end{equation*}
    \begin{equation*}
    \begin{aligned}
    \sum_{i=0}^nF_i = F_{n+2} - 1
    \end{aligned}
    \end{equation*}
    \begin{equation*}
    \begin{aligned}
    \sum_{i=0}^n F_i^2 = F_n F_{n+1}
    \end{aligned}
    \end{equation*}
    \begin{equation*}
    \begin{aligned}
    F_n^2 - F_{n-1}F_{n+1} = (-1)^{n+1}
    \end{aligned}
    \end{equation*}
    \begin{equation*}
    \begin{aligned}
    \gcd(F_m, F_n) = F_{\gcd(m, n)}
    \end{aligned}
    \end{equation*} \\ 
    $1, 1, 2, 3, 5, 8, 13, 21, 34, 55, 89, 144, 233, 377, 610, 987$

    \item \textbf{Derangement}
\begin{minted}[fontsize=\small]{c++}
d[0] = 1, d[1] = 0;
for(int n=2; n<=N; n++) d[n] = (n-1) * (d[n-1] + d[n-2]) % MOD;
\end{minted}
    $!n$ = permutations of $n$ elements with no fixed points.

    \item \textbf{Harmonic Series} \\
    \begin{equation*}
    \begin{aligned}
    H_n = \sum_{k=1}^{n} \frac{1}{k} \approx \ln n + \gamma + \frac{1}{2n} - \frac{1}{12n^2} + \frac{1}{120n^4}
    \end{aligned}
    \end{equation*} \\ 
    where $\gamma \approx 0.5772156649$ is the Euler-Mascheroni constant.

    \item \textbf{Stirling's Approximation} \\ 
    \begin{equation*}
    \begin{aligned}
    n! \sim \sqrt{2\pi n}\left(\frac{n}{e}\right)^n
    \end{aligned}
    \end{equation*}

    \item \textbf{Taylor Series} \\
    \begin{equation*}
    \begin{gathered}
    e^x = \sum_{n=0}^{\infty} \frac{x^n}{n!}, \quad
\sin x = \sum_{n=0}^{\infty} (-1)^n \frac{x^{2n+1}}{(2n+1)!} \\
\cos x = \sum_{n=0}^{\infty} (-1)^n \frac{x^{2n}}{(2n)!}, \quad
\ln(1+x) = \sum_{n=1}^{\infty} (-1)^{n+1} \frac{x^n}{n}
    \end{gathered}
    \end{equation*}

    \item \textbf{Lifting the exponent lemma} \\
    \begin{equation*}
    \begin{aligned}
    v_p(a^n - b^n) &= v_p(a-b) + v_p(n)
\quad (p \mid a-b,\ p \nmid ab) \\
v_p(a^n - b^n) &= v_p(a-b) + v_p(a+b) + v_p(n) - 1
\quad (p = 2,\ n \text{ even})
    \end{aligned}
    \end{equation*}

%     \item \textbf{Matroid} \\
%     A matroid $M = (E, \mathcal{I})$ is a ground set $E$ with a family of ``independent'' subsets
%     $\mathcal{I} \subseteq 2^E$ satisfying:
%     (I1) $\emptyset \in \mathcal{I}$;
%     (I2) hereditary: $B \in \mathcal{I}, A \subseteq B \implies A \in \mathcal{I}$;
%     (I3) exchange: $A, B \in \mathcal{I}, |A| < |B| \implies \exists x \in B \setminus A$ with
%     $A \cup \{x\} \in \mathcal{I}$. \\
%     A maximal independent set is a \emph{basis}; by (I3) all bases have the same size, the
%     \emph{rank} $r(M)$. The greedy algorithm (sort elements by weight, add each if it keeps the
%     current set independent) finds a max-weight independent set / basis for any matroid, and only
%     for matroids (this is essentially matroid's defining property for optimization).
%     Common examples: \emph{graphic matroid} ($E$ = edges, independent = forests -- greedy on
%     this is Kruskal's MST), \emph{linear matroid} ($E$ = vectors, independent = linearly
%     independent sets), \emph{uniform matroid} $U_{k,n}$ (independent = any subset of size $\le k$),
%     \emph{transversal matroid} of a bipartite graph $G=(X,Y,E)$ ($A \subseteq X$ independent iff
%     matchable into $Y$ -- see \emph{Matching is a Linear Matroid} for its linear representation).
% 
%     \item \textbf{Matroid Intersection Theorem} \\
%     For two matroids $M_1=(E,\mathcal{I}_1)$, $M_2=(E,\mathcal{I}_2)$ on the same ground set,
%     \begin{equation*}
%     \begin{aligned}
%     \max_{I \in \mathcal{I}_1 \cap \mathcal{I}_2} |I| = \min_{A \subseteq E} \big(r_1(A) + r_2(E \setminus A)\big)
%     \end{aligned}
%     \end{equation*}
%     where $r_1, r_2$ are the rank functions of $M_1, M_2$. A maximum common independent set can be
%     found in polynomial time by repeatedly augmenting along a shortest path in the exchange graph
%     (arcs $y \to x$ for $x \in I, y \notin I$ if $I - x + y \in \mathcal{I}_1$, and $x \to y$ if
%     $I - x + y \in \mathcal{I}_2$, plus source arcs from $\{y : I+y \in \mathcal{I}_1\}$ and sink
%     arcs to $\{y : I+y \in \mathcal{I}_2\}$), analogous to augmenting paths in bipartite matching.
%     Bipartite matching itself is the intersection of two partition (transversal) matroids on the
%     edge set, one per side; this is why matching, arborescence, and colored-forest problems all
%     reduce to matroid intersection. (Three or more matroids: intersection is NP-hard in general.)

\end{itemize}